\clearpage
\section{第 8.2 节教材正文例题}
\label{exercise-set:QM-C08-S02}
\exercisemeta
  {QM-C08-S02-EX}
  {2026-09-24 至 2026-09-25}
  {Shankar，第 8.2 节正文例子；印刷页 pp.~224--225；PDF pp.~235--236；式 (8.2.1)--(8.2.2)，图 8.3}
  {教材未编号讲解例；恢复量纲与积分步骤为展开推导}
  {例题讲解已完成；笔记已确认}

\exercisequestion{QM-C08-S02-E01}{两条自由粒子路径的作用量与相位差}
\textbf{对应知识点：}\relatedtopic{QM-C08-S02-T02}{作用量差与相干范围}。

\textbf{例题。}质量为 $m$ 的自由粒子在 $t=0$ 时位于原点，在 $t=T$ 时到达 $x=\ell$，其中 $T,\ell>0$。比较两条连接相同端点的路径
\begin{equation}
  x_{\mathrm{cl}}(t)=\ell\frac{t}{T},\qquad
  x_{\mathrm{alt}}(t)=\ell\left(\frac{t}{T}\right)^2.
\end{equation}
计算两者的作用量及相位差，并在 $\ell=1\,\mathrm{cm}$、$T=1\,\mathrm{s}$ 时比较质量为 $1\,\mathrm g$ 的粒子与电子。教材以厘米、秒为单位把这两式简写为 $x=t$ 与 $x=t^2$；这里显式保留量纲。

\begin{checkedsolution}
\textbf{两条路径与运动方程。}二者均满足 $x(0)=0$、$x(T)=\ell$，但
\begin{equation}
  \dot x_{\mathrm{cl}}=\frac{\ell}{T},\quad
  \ddot x_{\mathrm{cl}}=0;\qquad
  \dot x_{\mathrm{alt}}=\frac{2\ell t}{T^2},\quad
  \ddot x_{\mathrm{alt}}=\frac{2\ell}{T^2}.
\end{equation}
只有前者满足自由粒子的经典运动方程。后者仍是路径积分的候选路径：固定端点不意味着每条路径都必须满足 Newton 方程。

\textbf{作用量积分。}自由粒子的 $V=0$，故 $L=\frac12m\dot x^2$。经典路径给出
\begin{equation}
  S_{\mathrm{cl}}
  =\int_0^T\frac12m\frac{\ell^2}{T^2}\dd t
  =\boxed{\frac{m\ell^2}{2T}}.
\end{equation}
另一条路径给出
\begin{align}
  S_{\mathrm{alt}}
  &=\int_0^T\frac12m\left(\frac{2\ell t}{T^2}\right)^2\dd t
    =\frac{2m\ell^2}{T^4}\int_0^Tt^2\dd t\\
  &=\frac{2m\ell^2}{T^4}\frac{T^3}{3}
    =\boxed{\frac{2m\ell^2}{3T}}.
\end{align}
所以
\begin{equation}
  \boxed{\Delta S=S_{\mathrm{alt}}-S_{\mathrm{cl}}
    =\frac{m\ell^2}{6T}},\qquad
  \boxed{\Delta\theta=\frac{\Delta S}{\hbar}
    =\frac{m\ell^2}{6\hbar T}}.
  \label{eq:QM-C08-S02-free-phase}
\end{equation}
其中 $m\ell^2/T$ 的单位为 $\mathrm{kg\,m^2/s}=\mathrm{J\,s}$，与 $\hbar$ 相同。

\noindent\begin{minipage}{\linewidth}
\textbf{宏观与微观比较。}取 $\ell=10^{-2}\,\mathrm m$、$T=1\,\mathrm s$，使用 $\hbar\simeq1.055\times10^{-34}\,\mathrm{J\,s}$：
\begin{center}
\begin{tabular}{@{}llll@{}}
  \toprule
  对象 & 质量 $m$ & $\Delta S\ (\mathrm{J\,s})$ & $\Delta\theta\ (\mathrm{rad})$\\
  \midrule
  宏观粒子 & $10^{-3}\,\mathrm{kg}$ & $1.67\times10^{-8}$ & $1.58\times10^{26}$\\
  电子 & $9.11\times10^{-31}\,\mathrm{kg}$ & $1.52\times10^{-35}$ & $0.144$\\
  \bottomrule
\end{tabular}
\end{center}
\end{minipage}\par
教材粗略取电子质量为 $10^{-27}\,\mathrm g$、$\hbar\simeq10^{-27}\,\mathrm{erg\,s}$，因而得到 $\Delta\theta\simeq1/6$；与表中的较精细估值表达同一数量级结论。

对宏观粒子，同样的路径变形对应极强的相位变化；在平稳相位近似下，相干贡献集中于经典路径的狭窄邻域。对电子，这条明显弯曲的候选路径与经典路径仍只相差约 $0.14$ rad，不能仅依据偏离经典轨迹就将其忽略。此比较展示相位变化的尺度，不是用两项振幅代替完整路径积分，也不是凭两个相位相差巨大就认定它们恰好抵消。
\end{checkedsolution}
